\section{Discussion}
\label{dist:sec:discussion}

The single-cell description is now complete in distribution: geometric states
at every time, a geometric quasi-stationary limit, and a burst-size law
identical to it, with the conditional laws of burst time and burst size, the
multi-founder extension and the chained-transfer model alongside. What follows
reads that description into measurement, into two published parameter sets,
and into what is still open; the population consequences are taken up in
\ChPop.

\subsection{Single-cell assay predictions for lytic systems}
\label{dist:sec:discussion-assay}

For a genuinely lytic system --- \yp{}, \ft{}, lytic phage --- a
limiting-dilution assay, one infected cell per well with counts followed in
time, measures the single-cell process directly, with no population model
interposed. The number of new infectious particles released per cell is an
integer-valued \rv, and its full distribution, not merely its mean,
determines whether a small inoculum grows into a detectable infection
\cite{williams2024reproduction,pearson2011stochastic}. The theory makes five
predictions, all parameter-free given $(\lambda,\mu,\delta)$:
\begin{enumerate}[leftmargin=*,itemsep=1pt]
\item a fraction $b$ of wells show no release ever --- the mass of internal
extinction;
\item the cumulative release in a well is a step function: zero, then a
single terminal jump;
\item the jump times follow the density $\delta J(t)/(1-b)$ of
\cref{dist:eq:phi};
\item the jump sizes follow the geometric law with ratio $1/a$ --- the
quasi-stationary distribution, by \cref{dist:thm:identity};
\item the between-well variance of total release is $\var{W_\infty}$ from
\cref{dist:eq:ew2}, a quantity no deterministic constant-rate model can
produce at all.
\end{enumerate}
The five are not all of a kind. The second tests a standing assumption rather
than a derived result --- instantaneous complete release is assumption 3 of
\cref{dist:sec:assumptions}, and an assay that saw graded release would
refute the premise before reaching the theory --- and the first is a
one-parameter consequence of the roots alone. Predictions 3 to 5 are where
the chapter's own content is exposed: a density, a distribution and a
variance, each fixed by $(\lambda,\mu,\delta)$ with nothing left to tune.
Stated as literally as that, the predictions are too strong for a budding
virus: intermittent budding is not a single terminal burst
\cite{hataye2019principles}. They belong to lytic hosts, or to the bursting
endpoint of the spectrum of release models \ChPop{} lays out.

\subsection{The same theory at two extremes}
\label{dist:sec:two-extremes}

The theory has already been fitted to two intracellular pathogens whose
parameter estimates sit at opposite ends of what the model allows, and the
contrast shows how much behaviour is fixed by $(\lambda,\mu,\delta)$ alone.
\Cref{dist:tab:twoextremes} evaluates the closed forms at both.

\begin{table}[tbp]
\centering
\renewcommand{\arraystretch}{1.2}
\small
\begin{tabular}{@{}lcc@{}}
\toprule
& \ft{} SCHU S4 & \textit{B.~anthracis} \\
\midrule
$(\lambda,\mu,\delta)$, h$^{-1}$ & $(0.15,\ 0.01,\ 1.5\times10^{-4})$ & $(0.64,\ 1.64,\ 0.04)$ \\
$a$ & $1.001$ & $2.66$ \\
$b$, never bursts & $6.66\,\%$ & $96.24\,\%$ \\
$\E{\Kb\mid\text{burst}}=a/(a-1)$ & $934$ & $1.60$ \\
median burst size & $648$ & $1$ \\
$\E{T_{\rm prod}}$ & $45.6$ h & $0.74$ h \\
$\E{\tau\mid\text{burst}}$ & $48.4$ h & $0.93$ h \\
\bottomrule
\end{tabular}
\caption{The same three-rate theory at two published parameter sets, from
\cite{carruthers2020stochastic} and \cite{williams2021anthrax}. Every entry is
a closed form of this chapter evaluated at those rates:
\cref{dist:eq:roots,dist:eq:vinf,dist:eq:qsd,dist:eq:tprod,dist:eq:tauburst}.
The digits shown are those of the evaluation, not of the estimate: each entry
is exact given the rates in the first row, and inherits whatever uncertainty
those rates carry. The two columns come from different methods applied to
different data, so the comparison between them is qualitative even where the
arithmetic within a column is not.}
\label{dist:tab:twoextremes}
\end{table}

For \ft{} strain SCHU S4, modelling of murine challenge data gave
$\lambda=0.15$, $\mu=0.01$ and $\delta=1.5\times10^{-4}$ per hour
\cite{carruthers2020stochastic}. At those rates fewer than $7\%$ of infected
cells fail to burst in the model, and those that rupture release a geometric
number of bacteria with mean $934$ and median $648$. Bursts of that size sit
comfortably with the observed ability of doses as low as ten colony-forming
units to establish infection \cite{oyston2004tularaemia}: on the model's
account each rupture floods the local environment heavily enough that a small
inoculum needs only a few rounds of the cycle to reach large numbers.

For \textit{Bacillus anthracis}, approximate Bayesian computation against
published time series produced very different estimates:
$\lambda\approx0.64$, $\mu\approx1.64$, $\delta\approx0.04$ per hour
\cite{williams2021anthrax}. Here death dominates birth, and the model has more
than $96\%$ of phagocytes clearing their infection without rupturing, the rest
releasing about $1.6$ bacteria on average. The reported infectious dose,
$8\times10^3$ to $5\times10^4$ spores \cite{williams2021anthrax}, is orders of
magnitude above the tularaemia figure, and matches a pathogen whose
intracellular dynamics destroy far more than they release.

Two rows of \cref{dist:tab:twoextremes} are directly falsifiable, and they are
the rows the chapter had not previously computed. On these estimates a
\ft{}-infected macrophage should rupture at around two days post-infection,
and an anthrax-infected phagocyte, in the rare event that it ruptures at all,
within the hour. Neither number has a measurement to sit beside it here. What
would settle them is time-lapse imaging of infected macrophages held to
rupture, with the time of the release event recorded and the payload counted
directly; the prediction is sharp enough that a handful of such recordings
would be decisive, and sharp enough to be wrong.
% TODO: if a time-lapse rupture study for F. tularensis macrophages can be
% cited, restore the stronger claim and reference it here.
Both timings are only as good as the rates they condition on, and those rates
came from different methods and different data. What the pair shows regardless
is one three-rate theory spanning the space from near-certain large bursts to
near-certain clearance with minimal release.

\subsection{What the two probes showed}
\label{dist:sec:probes}

The two probes of \cref{dist:sec:moi,dist:sec:chained} were asked the same
question --- how far does the geometric structure reach --- and returned
opposite answers, which is the more informative outcome.

Multiplicity leaves the structure intact. Lineages are independent until the
shared clock rings, so $G_k=G^k$ and the load in a multiply infected cell is a
sum of independent copies of the single-founder load; the bookkeeping changes,
and two natural guesses about it fail
(\cref{dist:rem:ihatk,dist:rem:gk}), but the geometric shape does not.
Chaining destroys it as soon as internal extinction is possible, because the
$\mu=0$ decomposition needs a cell to have exactly one way to stop and the
death rate gives it two. \emph{The geometric structure is robust to
multiplicity and fragile to internal extinction in the chained setting}, and
the two failures are failures of different things --- a
counting identity in one case, a decomposition of the process in the other.

\subsection{Open problems}
\label{dist:sec:open-problems}

\begin{enumerate}[leftmargin=*,itemsep=2pt]
\item \textbf{How far the telescoping criterion reaches.}
\Cref{dist:thm:identity} closes what was an open problem: the identity
between burst size and the quasi-stationary law has a structural proof, and
\cref{dist:rem:criterion} states the property that proof uses. Two questions
survive it. Does that property \emph{characterise} the processes for which
the identity holds, or is it merely sufficient? And can the identity be
obtained independently by a Yaglom-type argument
\cite{yaglom1947certain}, which would connect it to the existence theory
rather than to the calculus?
\item \textbf{Chained joint laws for general $\mu$.} The joint laws of
$(t_k,r_k)$ for $\mu>0$, and the joint law of $(t_M,r_M)$ for the full
$M$-cell chain, remain open. The coefficient machinery of
\cref{dist:app:extract} is the available route, and
\cref{dist:sec:chained-scope} says why the elementary route is not.
\item \textbf{Logistic intracellular growth.} The standing objection to linear
birth is that it lets the conditional load grow without bound. The defence is
that the burst hazard is proportional to load and so regulates it ---
\cref{dist:thm:qsd} is precisely the statement that the conditional load
stabilises at $a/(a-1)$ --- but the sensitivity of the burst statistics, and
of the kernels \ChPop{} will use, to a logistic birth rate remains to be
checked numerically.
\end{enumerate}

\subsection{What comes next}
\label{dist:sec:next}

Everything in this chapter is single-cell. The released particles have so far
been counted and then forgotten, and that is the approximation \ChPop{}
removes. There the survival and release statistics derived here become the
age-structured kernels of a population-level renewal model: cells infected at
different times contribute to the infected-cell and free-particle pools
according to $\Ifix(a)$ and $g(a)=\delta K(a)$, in the tradition of
age-structured renewal running from
\cite{mckendrick1926applications,vonfoerster1959some}. The constant-rate
models of viral dynamics then appear as a projection of that renewal system,
and bursting is compared against budding at the level where the comparison
matters --- establishment, invasion speed, and what data can and cannot
identify.

One notational matter is left for later. \fwd{twotype}{\ChTwoTypecap}
solves a Riccati equation structurally identical to \cref{dist:eq:riccati}
under unrelated names, as \cref{dist:sec:roots} noted; aligning the two would
mean renaming $a$ and $b$ throughout this chapter, \ChCore{} and \ChPop, and
is a thesis-wide decision rather than a local one.
